\documentclass[APA,Utopia2COL]{WileyNJDv5}

\usepackage{amsmath,amssymb,amsthm,booktabs}
\usepackage{graphicx}
\hypersetup{hidelinks}

\articletype{Research Article}
\journal{Journal of Forecasting}
\raggedbottom

\begin{document}

\title{Dynamic Forecast-Optimal Smoothness for Penalized Trend Estimation}

\author[1]{Heriberto Espino Montelongo}
\authormark{ESPINO MONTELONGO}
\titlemark{DYNAMIC FORECAST-OPTIMAL SMOOTHNESS}

\address[1]{\orgname{Universidad de las Am\'ericas Puebla}, \orgaddress{\state{Puebla}, \country{Mexico}}}
\corres{Corresponding author: Heriberto Espino Montelongo. \email{[corresponding email]}}

\abstract[Abstract]{
Finite-difference penalized least squares yields a family of trend estimators indexed by smoothness, but chronological forecast-loss surfaces can be multimodal and their local minima can move through time. This paper studies smoothness selection as a dynamic forecasting problem. We first retain pooled forecast cross-validation as a static baseline: the smoothness value minimizing average historical future-block forecast loss. We then recover the local minima of origin-specific forecast-loss surfaces, connect nearby minima across adjacent origins into data-driven branches, and record each branch through its smoothness path together with Validation-1 and subsequent Validation-2 forecast loss. The forecasting decision is separated into two steps: select a historically persistent branch, then map that branch history into the current smoothness using rules such as the newest minimum, recent mean or median, and loss- or recency-weighted summaries. After smoothness is chosen, the trend is always refit on the newest information available before forecasting the untouched future block. The framework preserves the normalized smoothness coordinate and effective-degrees-of-freedom interpretation of controlled-smoothness penalized trends while allowing temporal persistence in multiple forecast-optimal smoothing regimes to enter the decision. A companion numerical paper addresses recovery and temporal correspondence of the local minima themselves.
}

\keywords{forecast evaluation; penalized least squares; smoothing parameter; trend estimation; time-series cross-validation; rolling-origin forecasting}

\maketitle

\section{Introduction}

Smoothing parameters are often introduced as technical constants, although in forecasting they determine a substantive modeling choice: how much recent movement is treated as signal and how much is treated as noise. A small penalty lets the fitted trend track the observations closely, while a large penalty forces the fitted path toward a low-degree polynomial. If that trend is subsequently extrapolated, the smoothing choice also changes the future trajectory.

Finite-difference penalized least squares (PLS) makes this trade-off especially transparent. Guerrero \cite{guerrero2007penalized,guerrero2008smoothness} showed that the penalty can be translated into an interpretable percentage of smoothness. This controlled-smoothness view is useful because the same numerical penalty does not represent the same effective amount of smoothing across sample sizes and difference orders. The remaining practical question is immediate: if the fitted trend will be used for forecasting, which percentage of smoothness should be chosen?

The forecasting literature provides strong precedents for choosing tuning parameters by predictive performance. Forecast-error minimization has long been used to estimate exponential-smoothing parameters \cite{taylor2004stes}; time-series cross-validation selects a kernel bandwidth through one-step-ahead prediction in dependent data \cite{hart1994tscv}; and recent work in the \emph{Journal of Forecasting} continues to connect hyperparameter selection, trend filtering, and out-of-sample loss to forecast performance \cite{zafar2022filter,stanek2023oos,wolff2024stock,franjic2025predictor,xu2025vix}. These papers motivate forecast-oriented tuning, but they do not define the smoothness percentage of a finite-difference penalized trend through a horizon-specific rolling future-block loss.

For fixed difference order $d$, estimation-window length $L$, and forecast horizon $h$, let $S\in[0,1]$ denote normalized smoothness. At each historical forecast origin, the trend is fitted using only observations available at that origin, extrapolated by the continuation implied by the same finite-difference model, and scored against the next $h$ realized observations. Averaging these errors yields a chronological criterion $F_{d,L,h}(S)$. We define
\begin{equation}
\label{eq:main_target_intro}
S^\star_{d,L,h}
\in
\arg\min_{S\in[0,1]}F_{d,L,h}(S)
\end{equation}
as the \emph{forecast-optimal smoothness} for that forecasting configuration.

The proposal is deliberately narrower than a general automatic forecasting system. The order $d$, window length $L$, and horizon $h$ are treated as fixed when defining the continuous smoothness problem. Joint adaptation of all three quantities is a different model-selection problem. Likewise, the present paper defines and studies the forecast target; a companion numerical study treats the separate problem of efficiently recovering multiple minima when the one-dimensional objective is non-unimodal.

The distinction from existing smoothing criteria is also important. Classical cross-validation and generalized cross-validation are usually motivated by estimation or reconstruction performance, while information criteria optimize different statistical objectives. Cort\'es-Toto et al.\ \cite{cortes2017smoothness} study the smoothness achieved by several such selectors inside PLS. The criterion proposed here is not intended to dominate those methods universally. It answers a different question: conditional on using the penalized trend as a forecasting device with a specified continuation and horizon, which amount of smoothness minimizes chronological forecast loss?

The contribution has four parts. First, controlled smoothness is converted into an endogenous forecast-based selection rule. Second, the normalized $S\in[0,1]$ coordinate makes the unsmoothed and maximally smoothed finite-difference limits explicit. Third, the target is horizon-specific through genuine future blocks rather than only one-step reconstruction or leave-one-out error. Fourth, tuning is embedded in a strictly chronological rolling-origin protocol, so the information set matches actual forecast use.

The remainder of the paper positions the proposal relative to controlled smoothness, predictive smoothing selection, and forecast evaluation; develops the finite-difference PLS representation; defines the forecast-optimal criterion; establishes its basic properties; and specifies the simulation and empirical protocol required to evaluate when forecast-based smoothness differs materially from conventional alternatives.

\section{Related work and positioning}

\subsection{Controlled smoothness in finite-difference PLS}

The direct methodological foundation is the controlled-smoothness formulation of finite-difference PLS. Guerrero \cite{guerrero2007penalized} derives a statistical representation of the penalized trend estimator and introduces a trace-based smoothness index linking the penalty parameter to an interpretable amount of smoothness. Guerrero \cite{guerrero2008smoothness} develops this into a user-controlled trend-estimation procedure in which a desired smoothness level is chosen first and the corresponding penalty is recovered afterward.

Cort\'es-Toto et al.\ \cite{cortes2017smoothness} study what smoothness is produced when the PLS penalty is selected by CV, GCV, AICc, and BIC. Their results make an important distinction for the present paper: a rule for selecting $\lambda$ and the amount of smoothness ultimately induced by that rule are not the same object.

The present paper reverses the controlled-smoothness direction:
\[
\text{user chooses }S\longrightarrow\lambda(S)
\]
becomes
\[
\text{chronological forecast loss}\longrightarrow S^\star.
\]
The percentage of smoothness remains the interpretable object, but it is selected endogenously for the forecast task.

\subsection{Predictive selection of smoothing and tuning parameters}

Forecast-based tuning is established in other smoother families. Hart \cite{hart1994tscv} considers kernel smoothing with serially correlated errors and proposes time-series cross-validation (TSCV), selecting the bandwidth through one-step-ahead predictive performance. This is an important conceptual precedent because it distinguishes good trend recovery from good prediction under dependence.

The object considered here is different. Hart selects a kernel bandwidth and explicitly models serial dependence in the errors. We select the smoothness of a finite-difference PLS trend, retain the model's native continuation rule, and allow an arbitrary $h$-step future block to define the tuning target. The connection to Hart is therefore the predictive principle rather than estimator equivalence.

Forecast-error tuning is also standard in exponential smoothing. Taylor \cite{taylor2004stes} describes data-driven smoothing-parameter estimation through ex post one-step-ahead forecast errors and develops an adaptive smooth-transition formulation. More recent \emph{Journal of Forecasting} papers apply the same predictive principle to other tuning problems. Wolff and Echterling \cite{wolff2024stock} choose regularization and machine-learning tuning parameters using validation forecast errors. Franjic and Schweikert \cite{franjic2025predictor} propose a cross-validation strategy that connects predictor preselection directly to nowcast performance. Xu et al.\ \cite{xu2025vix} compare bandwidth-selection schemes in a time-varying-parameter HAR model and show that bandwidth choice affects forecast performance.

These papers establish that forecast-oriented tuning is a legitimate methodological target. The narrower gap addressed here is the selection of a controlled smoothness percentage inside the finite-difference PLS family.

\subsection{Trend filtering for forecasting}

Zafar et al.\ \cite{zafar2022filter} explicitly evaluate trend-filtering methods by the short-term forecasts they enable. Their study reinforces the idea that a trend representation should be judged partly by predictive consequences rather than only by visual smoothness or in-sample fit. Our objective is complementary: rather than introduce a new trend representation, we keep finite-difference PLS fixed and optimize the amount of smoothness within that family.

\subsection{Chronological forecast evaluation}

The proposed criterion uses pseudo-out-of-sample loss as a tuning objective. Tashman \cite{tashman2000oos} reviews out-of-sample forecast evaluation schemes, while Bergmeir and Ben\'itez \cite{bergmeir2012cv} discuss cross-validation for time-series predictors. Stan\v{e}k \cite{stanek2023oos} formalizes out-of-sample loss estimation and distinguishes rolling and fixed evaluation schemes. His formulation also makes explicit that the relevant forecast loss depends on the forecast horizon, model memory, and updating design.

This is closely aligned with $F_{d,L,h}(S)$. Smoothness is selected for a specified horizon and update protocol, and the information set available at each origin is part of the estimator's definition.

\section{Finite-difference penalized trend estimation}

\subsection{Penalized least squares}

At forecast origin $T$, let
\[
x_T=(y_{T-L+1},\ldots,y_T)^\top\in\mathbb{R}^{L}
\]
denote the estimation window. For difference order $d<L$, let
\[
D_d\in\mathbb{R}^{(L-d)\times L},
\qquad
Q=D_d^\top D_d.
\]
The baseline estimator solves
\begin{equation}
\label{eq:pls_problem}
\widehat\tau_T(\lambda)
=
\arg\min_{\tau\in\mathbb{R}^{L}}
\left\{
\|x_T-\tau\|_2^2+\lambda\|D_d\tau\|_2^2
\right\},
\qquad \lambda\ge0,
\end{equation}
with solution
\begin{equation}
\label{eq:H}
\widehat\tau_T(\lambda)
=
H_\lambda x_T,
\qquad
H_\lambda=(I+\lambda Q)^{-1}.
\end{equation}

The paper uses the zero-drift, identity-weight formulation in \eqref{eq:pls_problem} to isolate the smoothness-selection question. Guerrero's broader statistical formulation allows nonzero mean differences and more general covariance weighting \cite{guerrero2007penalized}; these are natural extensions but are not silently imposed in the baseline criterion.

\subsection{Spectral representation}

Because
\[
v^\top Qv=\|D_dv\|_2^2\ge0,
\]
$Q$ is symmetric positive semidefinite. Under the standard difference operator,
\[
\operatorname{rank}(Q)=L-d,
\]
so $Q$ has $d$ zero eigenvalues and $L-d$ positive eigenvalues. Write
\[
Q=
U\operatorname{diag}
\left(
\underbrace{0,\ldots,0}_{d},
\delta_1,\ldots,\delta_{L-d}
\right)U^\top,
\qquad
\delta_j>0.
\]
Then
\begin{equation}
\label{eq:H_spectral}
H_\lambda
=
U\operatorname{diag}
\left(
\underbrace{1,\ldots,1}_{d},
\frac{1}{1+\lambda\delta_1},
\ldots,
\frac{1}{1+\lambda\delta_{L-d}}
\right)U^\top.
\end{equation}

Thus the $d$ null-space directions are reproduced without shrinkage, while each penalizable direction is attenuated by
\[
\alpha_j(\lambda)=\frac{1}{1+\lambda\delta_j}.
\]
For $d=2$, the null space consists of discrete linear sequences, so increasing $\lambda$ moves the fitted path toward a linear trend.

\subsection{Normalized smoothness}

Guerrero's trace-based construction motivates measuring smoothing through the effective action of $H_\lambda$. We use
\begin{equation}
\label{eq:S}
S(\lambda)
=
1-
\frac{1}{L-d}
\sum_{j=1}^{L-d}
\frac{1}{1+\lambda\delta_j}.
\end{equation}

The normalization excludes the $d$ directions that cannot be penalized, so
\[
S(0)=0,
\qquad
\lim_{\lambda\to\infty}S(\lambda)=1.
\]

For a fixed linear smoother,
\[
\operatorname{edf}(\lambda)=\operatorname{tr}(H_\lambda).
\]
Combining \eqref{eq:H_spectral} and \eqref{eq:S},
\begin{equation}
\label{eq:edf}
\boxed{
\operatorname{edf}(\lambda)
=
L-(L-d)S(\lambda).
}
\end{equation}
Hence $S$ is a normalized complement of effective degrees of freedom.

\section{Forecast-optimal smoothness by chronological cross-validation}
\label{sec:forecast_cv}

\subsection{From a fitted trend to a prospective forecast}

The parameter of interest is the normalized smoothness $S$, not the
unbounded penalty $\lambda$. Define the continuous family of
smoothing matrices, including their limiting endpoint, by
\begin{equation}
\label{eq:sindexed_smoothing_matrix}
H(S)=
\begin{cases}
\bigl(I+\lambda(S)D_d^\top D_d\bigr)^{-1},&0\le S<1,\\
P_{\ker(D_d)},&S=1,
\end{cases}
\end{equation}
where $P_{\ker(D_d)}$ is the orthogonal projection onto the
null space of the finite-difference operator. Computing
$H(S)$ may require inverting the monotone map $S(\lambda)$,
but the domain on which cross-validation operates is the
\emph{full bounded interval} $[0,1]$.

A fitted trend is not automatically a forecast. We specify
a continuation rule that imposes zero future $d$th
differences, producing the fixed linear operator
\begin{equation}
\label{eq:continuation_operator_space}
G_{d,h}\in\mathbb{R}^{h\times L}.
\end{equation}
The corresponding prospective forecast from origin $t$ is
\begin{equation}
\label{eq:forecast_operator}
\widehat z_t(S)=G_{d,h}H(S)x_t.
\end{equation}
In terms of the final backward differences of the
historically fitted trend,
\begin{equation}
\label{eq:continuation}
\widehat\tau_{t+k\mid t}(S)
=\sum_{j=0}^{d-1}\binom{k+j-1}{j}
\nabla^j\bigl[H(S)x_t\bigr]_{L},
\qquad k=1,\ldots,h.
\end{equation}
For $d=1$, $d=2$, and $d=3$, these continuations are constant,
linear, and quadratic in the forecast step $k$, respectively.
More generally, fixing $d$ fixes the maximum degree $d-1$,
while changing $S$ changes the smoothed endpoint values and
differences, hence the coefficients of the forecast polynomial.
The method selects among \emph{future extrapolations}, not
merely among different post-hoc representations of the past.
Selecting $d$ itself is a separate model-selection decision;
it is held fixed within the baseline CV comparison.

\subsection{The forecast-MSE surface in the smoothness coordinate}

For an \emph{historical} forecast origin $t$, let
\begin{equation}
\label{eq:future_observation_block}
z_t=(y_{t+1},\ldots,y_{t+h})^\top
\end{equation}
be the subsequent realized validation block. The loss assigned
to a candidate smoothness $S\in[0,1]$ is
\begin{equation}
\label{eq:origin_loss_smoothness}
F_t(S)
=\frac1h\left\|z_t-G_{d,h}H(S)x_t\right\|_2^2.
\end{equation}
Expanding this expression makes the dependence on the
smoothness-indexed smoothing matrix explicit:
\begin{equation}
\label{eq:origin_loss_quadratic}
\begin{aligned}
h F_t(S)
={}&z_t^\top z_t
-2z_t^\top G_{d,h}H(S)x_t\\
&+x_t^\top H(S)G_{d,h}^\top
G_{d,h}H(S)x_t.
\end{aligned}
\end{equation}
The matrix $H(S)$ is symmetric. Unlike ordinary fitting
criteria, \cref{eq:origin_loss_smoothness} scores observations
\emph{after} the training endpoint; the in-sample residual
$\|x_t-H(S)x_t\|^2$ does not appear in the criterion.

For derivative calculations it is sometimes convenient to
use the algebraically equivalent auxiliary function
\begin{equation}
\label{eq:origin_loss}
f_t(\lambda)=F_t(S(\lambda)).
\end{equation}
This equality is a coordinate change, not a second
optimization problem or a requirement to search a truncated
range of $\lambda$.

\subsection{Chronological, horizon-matched forecast cross-validation}

At final forecast origin $T$, select historical origins
$t_1,\ldots,t_M$ whose training windows are available and
whose validation blocks are already observed:
\begin{equation}
\label{eq:chronological_origin_restriction}
t_m+h\le T,\qquad m=1,\ldots,M.
\end{equation}
For fixed $(d,L,h)$, the proposed CV surface is
\begin{equation}
\label{eq:F_S}
F^{\mathrm{pool}}_{T,d,L,h}(S)
=\frac1M\sum_{m=1}^M F_{t_m}(S),
\qquad S\in[0,1].
\end{equation}
It produces the forecast-optimal smoothness estimate
\begin{equation}
\label{eq:Sstar}
\widehat S^{\mathrm{FCV}}_{T,d,L,h}
\in\operatorname*{arg\,min}_{S\in[0,1]}
F^{\mathrm{pool}}_{T,d,L,h}(S).
\end{equation}
The selection space includes both exact limiting models,
$S=0$ and $S=1$, not just interior grid values.
When the objective has several minima, the global
minimizer is determined by the lowest historical predictive
MSE including the endpoints; a deterministic smallest-$S$
tie convention is used if several minima have the same
recorded objective value.

Both the forecasting horizon and the continuation operator
are essential parts of the criterion. Minimizing one-step
forecast error and deploying the resulting smoothness for
a longer horizon is a different procedure. The index makes
this comparison interpretable on $[0,1]$, while its
one-to-one transformation from $\lambda$ preserves the
fitted-trend family and any exact minimizing set.

\subsection{Final refit and untouched future forecast}

At outer origin $T$, all validation-stage trends are
discarded. Using the smoothness chosen from historical
blocks, the current trend is refitted on the newest
$L$ observations:
\begin{equation}
\label{eq:final_refit}
\widehat\tau_T
=H\bigl(\widehat S^{\mathrm{FCV}}_{T,d,L,h}\bigr)
y_{T-L+1:T}.
\end{equation}
The operational forecast is
\begin{equation}
\label{eq:final_operational_forecast}
\widehat y_{T+1:T+h\mid T}
=G_{d,h}\widehat\tau_T.
\end{equation}
No observation after $T$ enters either the CV decision
or the final fit. Historical forecast losses, not the
historical fitted trend paths, are averaged.

\subsection{Relationship to established cross-validation}

Ordinary leave-one-out CV and GCV primarily measure
interpolation or fitted-signal error for a smoother.
The present criterion uses chronological, genuinely
future-block MSE from a declared polynomial continuation.
This does not invent the general principle of predictive
cross-validation: Hart's time-series CV
\cite{hart1994tscv} is an important precedent.
The specific object investigated here is direct
$h$-step forecast-MSE selection for the
finite-difference PLS trend over normalized $S\in[0,1]$.
Time-adaptive decisions using historical minima, considered
later, are optional and not part of this definition.

\section{Basic properties}

\begin{proposition}[Monotone smoothness reparameterization]
\label{prop:monotone}
For positive penalized eigenvalues $\delta_1,\ldots,\delta_{L-d}$,
\begin{equation}
\label{eq:smoothness_first_derivative}
S'(\lambda)
=
\frac{1}{L-d}
\sum_{j=1}^{L-d}
\frac{\delta_j}{(1+\lambda\delta_j)^2}
>0,
\end{equation}
and
\begin{equation}
\label{eq:smoothness_second_derivative}
S''(\lambda)
=
-\frac{2}{L-d}
\sum_{j=1}^{L-d}
\frac{\delta_j^2}{(1+\lambda\delta_j)^3}
<0.
\end{equation}
After adjoining $\lambda=\infty$, $S$ maps $[0,\infty]$ continuously and strictly increasingly onto $[0,1]$.
\end{proposition}

\begin{proof}
Differentiate \cref{eq:S} term by term. Positivity of each $\delta_j$ gives the signs; the endpoint values follow from $(1+\lambda\delta_j)^{-1}$.
\end{proof}

The bounded scale has a common interpretation across fixed $(d,L)$ configurations while preserving the ordering of penalty strength.

\begin{proposition}[Effective degrees of freedom]
\label{prop:edf}
For fixed $\lambda$,
\begin{equation}
\label{eq:edf_smoothness_formula}
\operatorname{tr}(H_\lambda)=L-(L-d)S(\lambda).
\end{equation}
\end{proposition}

\begin{proof}
From \cref{eq:H_spectral},
\begin{equation}
\label{eq:edf_spectral_trace}
\operatorname{tr}(H_\lambda)
=
d+\sum_{j=1}^{L-d}(1+\lambda\delta_j)^{-1}.
\end{equation}
Substitution of \cref{eq:S} yields the result.
\end{proof}

\begin{proposition}[Equivalent optimization in penalty and smoothness space]
\label{prop:equiv}
Let $F(S)=f(\lambda(S))$. Then the global minimizers in $\lambda\in[0,\infty]$ and $S\in[0,1]$ correspond under the bijection in Proposition \cref{prop:monotone}. At interior points,
\begin{equation}
\label{eq:loss_change_of_variables}
F'(S)=\frac{f'(\lambda)}{S'(\lambda)},
\end{equation}
so interior stationary points correspond exactly.
\end{proposition}

\begin{proof}
A strictly increasing bijection preserves the minimizing set. The derivative identity follows from the chain rule.
\end{proof}

\begin{proposition}[Exact endpoints and the unique limiting fit]
\label{prop:endpoints}
Under the conditions of \cref{prop:pls_wellposed}, $H(0)=I$.
Let $U_0\in\mathbb R^{L\times d}$ have orthonormal columns spanning
$\ker(D_d)$. At $S=1$, interpreted as $\lambda\to\infty$,
\begin{equation}
\label{eq:exact_limiting_projection}
H(1)=H_\infty=U_0U_0^\top=P_{\ker(D_d)}.
\end{equation}
Although $H_\infty$ is singular with rank $d$ and nullity $L-d$,
the limiting trend is uniquely defined by
\begin{equation}
\label{eq:exact_limiting_constrained_fit}
H_\infty x_T
=\operatorname*{arg\,min}_{\tau:\,D_d\tau=0}
\|x_T-\tau\|_2^2.
\end{equation}
Thus no $1-\varepsilon$ surrogate is required for $S=1$.
\end{proposition}

\begin{proof}
At $\lambda=0$, $H_0=I$. By \cref{eq:H_spectral}, the
$d$ unit eigenvalues on $\ker(D_d)$ remain one, whereas
$(1+\lambda\delta_j)^{-1}\to0$ for every $\delta_j>0$.
This proves the projection formula. Any feasible
$\tau\in\ker(D_d)$ has the form $\tau=U_0a$.
Since $U_0^\top U_0=I$, the least-squares objective
$\|x_T-U_0a\|_2^2$ has the unique minimizer
$a=U_0^\top x_T$. Hence $\tau=U_0U_0^\top x_T$ is unique,
even though the projection matrix itself is not invertible.
\end{proof}

For $d=2$, the endpoints are the unsmoothed historical path and
its unique least-squares linear projection. The expression
$(I+\infty Q)^{-1}$ is not an ordinary matrix inverse:
$S=1$ must be computed as the exact spectral projection.
A large finite penalty or $S=1-\varepsilon$ is only an
approximation and can introduce avoidable numerical error.
Endpoint comparison is part of the scientific selection problem,
not merely a numerical detail.

\subsection{The future polynomial family}

The estimator produces a polynomial continuation only after the
future-difference rule has been specified. This matters for
interpreting forecast-CV as a prospective trend-selection method.

\begin{proposition}[Forecast polynomial indexed by smoothness]
\label{prop:continuation_polynomial}
Fix $1\le d<L$ and $h\ge1$. For every $S\in[0,1]$, the
continuation $k\mapsto\widehat\tau_{T+k\mid T}(S)$ is the
restriction to positive integer steps of a polynomial in
$k$ of degree at most $d-1$. The coefficients depend on
$S$ through the terminal backward differences of $H(S)x_T$.
\end{proposition}

\begin{proof}
For a fitted sequence $u=H(S)x_T$, imposing zero future
$d$th differences gives the discrete Newton continuation
\begin{equation}
\label{eq:newton_forecast_polynomial}
\widehat\tau_{T+k\mid T}(S)
=\sum_{j=0}^{d-1}\binom{k+j-1}{j}
(\nabla^j u)_L.
\end{equation}
Each binomial coefficient is a polynomial in $k$ of
degree $j$. The sum therefore has degree at most
$d-1$. Since $u=H(S)x_T$, its terminal differences
and consequently its coefficients vary with $S$.
The limiting case $S=1$ is well defined by the
orthogonal projection in \cref{prop:endpoints}.
\end{proof}

This proposition does not imply that the degree $d-1$
is chosen by forecast-CV: the baseline fixes $d$ and
selects among the corresponding extrapolations by
their historical future-block error. Changing $d$
would require an explicitly specified joint or
nested model-selection experiment.

\subsection{Existence and possible nonuniqueness}

The fitting problem has a unique solution for every $S\in[0,1]$: for $S<1$ this follows from strict convexity in \cref{prop:pls_wellposed}, and for $S=1$ from the constrained least-squares argument in \cref{prop:endpoints}. For finite $S<1$, $H_{\lambda(S)}$ depends continuously on $S$; the exact projection limit at $S=1$ makes the full family continuous on $[0,1]$. Thus, with a fixed finite collection of historical folds, $F^{\mathrm{pool}}_{T,d,L,h}(S)$ is continuous on the compact interval $[0,1]$ and attains at least one global minimum. \emph{Existence does not imply a unique minimizing $S$.}

Uniqueness is not assumed. Distinct smoothing regimes can produce similar validation loss, and the criterion may contain several local minima. This is statistically meaningful because competing minima can imply different effective degrees of freedom and terminal slopes. The implemented numerical search uses adaptive derivative bracketing, Brent root refinement, and exact endpoint comparison. The mathematical criterion in \cref{eq:Sstar} is independent of numerical search details and does not require tracking minima through time.

\subsection{Forecast optimality versus recovery optimality}

If a latent trend $\tau$ is known in simulation, a recovery-oriented oracle might choose
\begin{equation}
\label{eq:latent_recovery_oracle}
S_{\mathrm{rec}}
\in
\arg\min_S
\|\widehat\tau_S-\tau\|^2.
\end{equation}
The proposed criterion instead chooses $S^\star_{d,L,h}$ from future predictive error. These targets need not coincide. A smoother can reconstruct the historical latent path accurately but extrapolate poorly, or sacrifice historical recovery in exchange for more stable terminal derivatives.

This distinction is a central empirical question of the paper, not a claim that forecast-optimal smoothness is universally preferable under every loss.

\section{Evaluation protocol}

The empirical program now has two layers. Checkpoint 03 is a frozen paper-scale evaluation of the simpler pooled forecast-CV selector and remains a baseline. The central dynamic extension evaluates tracked local-minimum branches and alternative rules for converting branch history into the smoothness used at the current forecast origin.

\subsection{Layer 1: pooled forecast-CV baseline}

The frozen pooled baseline fixes
\[
d=2,\qquad L=60,\qquad h\in\{1,3,6,12\}.
\]
It compares horizon-matched pooled forecast-CV with one-step forecast-CV, ordinary CV, GCV, and AICc. BIC remains an exploratory diagnostic because its globally searched score was pinned to the left finite boundary in the refinement stage.

The pooled study uses controlled latent mechanisms, iid Gaussian, AR(1), and finite-variance Student-$t$ noise, fixed outer origins, and untouched future blocks. Its design was frozen before the paper-scale Checkpoint 03 run and is not retuned after observing those results.

\subsection{Layer 2: dynamic tracked-minimum evaluation}

The central dynamic experiment operates on a chronological sequence of Validation-1 forecast-loss surfaces. At each origin $t$, all relevant local minima
\[
\mathcal M_t=\{S_{1,t},\ldots,S_{K_t,t}\}
\]
are recovered. Local minima are linked across adjacent origins by one-to-one matching under
\[
|S_{j,t}-S_{j,t-1}|\le\varepsilon.
\]
The resulting branch history is stored as
\[
V_j
=
\{(S_{j,t},\ell^{(1)}_{j,t},\ell^{(2)}_{j,t})\}_{t\in\mathcal T_j},
\]
where $\ell^{(1)}$ is Validation-1 forecast loss and $\ell^{(2)}$ is Validation-2 forecast loss after refitting through Validation 1.

The dynamic protocol separates branch selection from final smoothness selection. A branch rule
\[
\widehat j_T=\psi(V_1,\ldots,V_J)
\]
uses historical branch support and Validation-2 performance only. Conditional on the selected branch, a rule
\[
\widehat S_T=\phi(V_{\widehat j_T})
\]
produces the smoothness used for the current forecast.

The first comparison set for $\phi$ contains:
\begin{enumerate}
\item the newest tracked local minimum;
\item a recent arithmetic mean;
\item a recent median;
\item a Validation-2-loss-weighted mean;
\item a recency-and-Validation-2 weighted mean.
\end{enumerate}
A direct forecast of the smoothness trajectory is reserved as a later extension unless its specification can be frozen without using outer-test information.

\subsection{Selection, refitting, and information boundaries}

Validation selects hyperparameters or branch summaries; it does not select a historical trend fit to carry into the test. For every candidate branch rule, after $\widehat S_T$ is obtained, all temporary Validation-1 and Validation-2 fits are discarded. The trend is freshly estimated using the most recent full window available at the outer origin:
\[
x_T=y_{T-L+1:T},
\qquad
\widehat\tau_T
=
H_{\lambda(\widehat S_T)}x_T.
\]
The operational forecast is
\[
\widehat y_{T+1:T+h\mid T}
=
G_{d,h}\widehat\tau_T.
\]
The outer future block is revealed only for scoring. At a later outer origin, observations that have become available enter the historical information set and the complete branch-selection, smoothness-selection, refit, and forecast cycle is repeated.

Thus the procedure may average historical forecast losses or historical smoothness values when a declared rule requires it, but it never averages historical trend fits.

\subsection{Fair comparison of branch-to-smoothness rules}

All $\phi$ rules must be evaluated on exactly the same recovered branch histories, the same branch selector $\psi$, and the same untouched outer test blocks. Hyperparameters such as the continuation radius $\varepsilon$, recent-history length $K$, recency factor $\rho$, and weighting stabilizer $\delta$ must be fixed from development data before final test scoring.

The dynamic experiment should report forecast MSFE and MAE, selected smoothness, branch support, branch switching or disappearance rates, and the distance between each feasible rule and a simulation-only latent forecast oracle where available. The pooled forecast-CV selector remains a direct baseline.

\subsection{Role of the numerical-methods companion paper}

The forecasting paper assumes that the relevant local minima on each surface and their cross-origin correspondences can be supplied reliably. Numerical discovery of those minima, endpoint handling, adaptive subdivision, derivative-root refinement, and the branch-correspondence problem belong to the companion numerical-methods paper. The forecasting paper evaluates what information in the resulting branch matrices is useful for selecting the current smoothness.

\subsection{Public real-data evaluation}

The final paper should complement simulation with a heterogeneous public panel under the same chronology. No series-specific branch rule may be chosen from its test performance. If macroeconomic histories are used, historical backtests require real-time vintages rather than revised current histories.

\subsection{Freeze rules}

The frozen Checkpoint 03 pooled design remains unchanged. The dynamic branch experiment receives its own freeze point: the branch matcher, $\psi$, candidate $\phi$ rules, $\varepsilon$, $K$, $\rho$, $\delta$, horizons, origins, and primary metrics must be recorded before the final dynamic paper-scale run. Unfavorable final results do not justify changing those choices.

\section{Scope, interpretation, and implications}

\subsection{What the dynamic smoothness object means}

The selected smoothness is not treated as a fixed property of a series. At a given origin, several local minima of the forecast-loss surface may coexist. Their temporal branches summarize recurring smoothing regimes discovered by the data.

For a final selected value $\widehat S_T$, the effective-degrees-of-freedom interpretation remains
\begin{equation}
\label{eq:scope_smoothness_edf}
\widehat S_T
=
\frac{L-\operatorname{edf}_T}{L-d}.
\end{equation}
A value near zero corresponds to a flexible trend; a value near one corresponds to a trend close to the null-space polynomial. The dynamic contribution is that today's value can depend on the history of a persistent local-minimum branch rather than on one pooled optimum alone.

\subsection{Branches are data-driven, not fixed smoothness bins}

The continuation radius $\varepsilon$ does not define a preassigned category such as ``smoothness around 0.9.'' Branch locations arise from local minima of successive forecast-loss surfaces. The radius only determines whether two neighboring minima are close enough to be interpreted as the same branch.

Consequently, a branch can drift over time. Its identity reflects local continuity of the optimum, not membership in a fixed interval.

\subsection{What may be averaged}

The framework distinguishes three objects that should not be conflated. Historical forecast losses may be averaged to rank candidate smoothness values or branches. Historical smoothness values may be averaged when a declared branch functional such as a recent-mean rule is used. Historical fitted trends are never averaged in the proposed procedure.

After the current smoothness is selected, the trend is always refit on the newest information available before the forecast.

\subsection{Static and dynamic targets}

The pooled selector asks which single value of $S$ performs best on average over historical forecast origins. The branch formulation asks how a time-indexed collection of local minima can be organized and converted into multiple admissible current values of $S$. Predictive improvement is a possible outcome of particular decisions, not a property guaranteed by the representation itself.

Neither target is automatically the smoothness that best reconstructs a latent historical trend. Forecasting and signal recovery remain distinct objectives.

\subsection{Numerical boundary}

The forecasting paper does not claim that a particular root-search or assignment algorithm is the substantive contribution. Recovery of multiple minima, exact endpoint handling, and temporal correspondence of minima belong to the companion numerical-methods paper.


\subsection{Boundary behavior of smoothness forecasts}

Normalized smoothness is restricted to $[0,1]$. Convex means of historically
admissible values remain in that interval, whereas linear or increment
extrapolations of the branch trajectory need not. Truncation back to the
admissible interval is therefore part of defining these functionals,
not an innocuous plotting convention. The simulations show that the
frequency of truncation varies substantially with the form of $\phi$.
This is a concrete difference between equally admissible constructions
of the next smoothness value.

\subsection{Limits of predictive interpretation}

A favorable four-series confirmation result is not sufficient to establish
universal superiority. The broad financial panel and controlled changing-
roughness experiments did not reproduce an aggregate advantage over
pooled forecast-CV. Their errors are evaluated on different scales, and
the series within the financial panel are not independent.
Accordingly, the paper claims a reusable representation for chronologically
tracked forecast-loss minima and an illustrative family of decision maps;
it does not claim that tracked branches inevitably improve forecasting.

\subsection{Limitations}

The baseline estimator assumes equally spaced observations, a fixed difference order, a fixed final estimation window, and deterministic continuation of the estimated trend. Checkpoints 05 and 06 show that high-order native polynomial continuation can amplify small differences in selected $S$ into extreme forecast errors. In particular, restricting higher difference orders reduced the catastrophic error tail in a post-hoc diagnostic, but does not validate a revised forecasting selector. The current branch construction also introduces design quantities such as the tracking radius, recent-history length, and weighting parameters. Those quantities must be selected on development data and frozen before outer-test evaluation.

Local-minimum identity can become ambiguous when branches approach or cross. The initial greedy one-to-one matcher is therefore a transparent baseline rather than a mathematical definition of true branch identity. Sensitivity to the numerical correspondence rule must be reported separately from forecasting performance.

Finally, the framework isolates trend extrapolation. It does not yet add a separate stochastic residual forecast model. That restriction keeps the interpretation of any forecast gain tied to smoothness selection rather than to another predictive component.

\section{Conclusion}

Selecting smoothness for a trend that
will be extrapolated is a forecasting
problem, not merely a post-hoc
descriptive smoothing problem.
We formulated chronological,
horizon-matched forecast
cross-validation over a normalized
PLS smoothness index $S\in[0,1]$.
The criterion compares historical
future-block errors from an explicitly
defined continuation, includes the
two limiting smoothing regimes, and
refits the selected trend before
issuing a genuinely out-of-sample
forecast.

The PLS smoother, trace-based
smoothness index, and predictive
cross-validation principle are
established methods. The contribution
studied here is their particular
operational combination for
finite-difference trend extrapolation
and a frozen simulation comparison
with conventional smoothing criteria.
At horizons $3$, $6$, and $12$,
the simulated geometric RMSFE
ratios to one-step tuning were
$0.947$, $0.861$, and $0.762$,
respectively, supporting the value
of selecting smoothness for the
forecast horizon within this design.

The objective can have multiple
local minima, and changing from
$\lambda$ to $S$ does not remove them.
Consequently, practical selection
must compare competing numerical
candidates and both endpoints.
The evidence does not establish
universal superiority of forecast-CV,
optimality across all polynomial
continuation orders, or a general
gain from tracked-minimum adaptation.
The principal result is a transparent
prediction-targeted tuning criterion
whose modeling assumptions,
computational search, and
information set can be stated
and evaluated explicitly.


\bmsection*{Data availability}
The methodological development does not depend on proprietary data. The submission-stage empirical evaluation is designed around public forecasting data and reproducible real-time or archived series when historical vintages matter. Any exploratory current-vintage macroeconomic results are treated separately from the submission-grade real-time-vintage analysis.

\bmsection*{Code availability}
Reusable implementations of finite-difference penalized trend estimation, chronological forecast objectives, rolling-origin validation, and smoothness transformations are maintained in the \texttt{trend\_estimation} package at \url{https://github.com/heri-espino/Trend-Estimation}. The repository also contains the manuscript source, experiment protocols, and frozen result artifacts used for the paper.

\bmsection*{Conflict of interest}
The author declares no conflict of interest.

\bibliography{references}

\end{document}
